\section{Reductions that hold for every variety}\label{sec:reductions}

\subsection{Correspondences}
Let $B$ and $X$ be varieties, possibly disconnected, with $B$ of pure
dimension $b$. For an algebraic cycle $\Gamma$ of codimension $c$ on
$B\times X$ with rational coefficients, the map
\[
  \Gamma_{*}\colon H^{k}(B,\QQ)\to H^{k+2c-2b}(X,\QQ),\qquad
  y\mapsto\pr_{X*}\bigl(\cl(\Gamma)\cup\pr_{B}^{*}y\bigr),
\]
is a morphism of Hodge structures $H^{k}(B)(b-c)\to H^{k+2c-2b}(X)$. It
carries algebraic classes to algebraic classes, the composite of two such
maps is the map of the composite correspondence, and for a morphism
$f\colon B\to X$ the graph $\Gamma_{f}$ and its transpose give $f_{*}$ and
$f^{*}$ \cite[Chapter~16]{Ful98}. For a cycle with components of several
codimensions we add the maps of the components.

A pure rational Hodge structure $V$ of weight $w$ is \emph{polarizable} if
it carries a bilinear form $Q$ such that $Q(Cx,\bar x)>0$ for $x\ne0$, where
$C$ acts on $V^{a,b}$ by $i^{a-b}$. The cohomology of a variety and its Tate
twists are polarizable, by the Hodge--Riemann relations
\cite[Chapter~6]{Voi02b}. Polarizable Hodge structures are
semisimple.\footnote{If $W\subseteq V$ is a sub-Hodge structure, the
restriction of $Q$ to $W$ is again a polarization, hence nondegenerate, so
$V=W\oplus W^{\perp}$. The orthogonal $W^{\perp}$ is a sub-Hodge structure,
because the decomposition $V_{\CC}=\bigoplus V^{a,b}$ is orthogonal for the
hermitian form $Q(Cx,\bar y)$ and $W_{\CC}$ is the sum of its intersections
with the $V^{a,b}$.}

\begin{lemma}\label{lem:lift}
Let $\varphi\colon V\to W$ be a surjective morphism of polarizable Hodge
structures of weight $2p$. Then $\varphi$ maps the Hodge classes of $V$, the
rational classes of type $(p,p)$, onto those of $W$.
\end{lemma}

\begin{proof}
By semisimplicity the kernel of $\varphi$ has a complement $V'$ that is a
sub-Hodge structure, and $\varphi$ restricts to an isomorphism of Hodge
structures $V'\to W$. Its inverse carries a Hodge class of $W$ to a Hodge
class of $V'\subseteq V$.
\end{proof}

\begin{lemma}[Dominated cohomology]\label{lem:dominated}
Let $X$ be a variety, and let $B_{1},\dots,B_{r}$ be varieties with
algebraic cycles $\Gamma_{j}$ on $B_{j}\times X$, such that the images of
the maps $\Gamma_{j*}$ span $H^{2p}(X,\QQ)$. Then every Hodge class of degree
$2p$ on $X$ is a sum $\sum_{j}\Gamma_{j*}\gamma_{j}$ with $\gamma_{j}$ Hodge
classes on $B_{j}$. In particular, if the Hodge conjecture holds for every
$B_{j}$, it holds for $X$ in degree $2p$.
\end{lemma}

\begin{proof}
Write $\Gamma_{j}=\sum_{c}\Gamma_{j,c}$ with $\Gamma_{j,c}$ of pure
codimension $c$, and $b_{j}=\dim B_{j}$ (if $B_{j}$ is not equidimensional,
treat its components separately). The component $\Gamma_{j,c}$ maps into
degree $2p$ exactly from $H^{k}(B_{j})$ with $k=2p+2b_{j}-2c$, which is
even. The sum of the maps
\[
  \bigoplus_{j,c}H^{2p+2b_{j}-2c}(B_{j},\QQ)(b_{j}-c)\longrightarrow
  H^{2p}(X,\QQ)
\]
is a surjective morphism of polarizable Hodge structures of weight $2p$.
A Hodge class in $H^{2p+2b_{j}-2c}(B_{j})(b_{j}-c)$ is a Hodge class of
$B_{j}$. Apply \Cref{lem:lift}. Algebraic classes go to algebraic classes
under correspondences, which gives the last assertion.
\end{proof}

\subsection{Degree two and hard Lefschetz}

\begin{theorem}[Lefschetz \cite{Lef24}, Kodaira--Spencer \cite{KS53}]\label{thm:lef11}
Every Hodge class of degree two is algebraic.
\end{theorem}

\begin{proof}
The exponential sequence gives an exact sequence
\[
  H^{1}(X,\cO_{X}^{\times})\to H^{2}(X,\ZZ)\to H^{2}(X,\cO_{X}),
\]
whose second
map is the composite of $H^{2}(X,\ZZ)\to H^{2}(X,\CC)$ with the projection
to $H^{0,2}(X)$. An integral class of type $(1,1)$ therefore is the first
Chern class of a holomorphic line bundle, which is algebraic by GAGA, and
its first Chern class is the class of a divisor. A rational Hodge class has
an integral multiple.
\end{proof}

Let $h$ be the class of a hyperplane section of $X$, of dimension $n$, and
$L$ the cup product with $h$. By the hard Lefschetz theorem
$L^{n-k}\colon H^{k}(X,\QQ)\to H^{2n-k}(X,\QQ)$ is an isomorphism for
$k\le n$ \cite[Theorem~6.25]{Voi02b}.

\begin{proposition}\label{prop:hardlef}
For $2p<n$, if every Hodge class of degree $2p$ on $X$ is algebraic, so is
every Hodge class of degree $2n-2p$. Every Hodge class of degree $0$, $2$,
$2n-2$ or $2n$ is algebraic, and \HC{} holds for every variety of dimension
at most three.
\end{proposition}

\begin{proof}
$L^{n-2p}$ maps $H^{p,p}$ onto $H^{n-p,n-p}$ and rational classes to rational
classes, so it maps $\Hdg^{p}(X)$ onto $\Hdg^{n-p}(X)$; and it maps the class
of a subvariety $Z$ to the class of $Z\cap H_{1}\cap\dots\cap H_{n-2p}$ for
general hyperplanes $H_{i}$. Hence
$\Hdg^{n-p}(X)=L^{n-2p}\Alg^{p}(X)\subseteq\Alg^{n-p}(X)$. Degrees $0$ and
$2n$ are spanned by the fundamental class and the class of a point; degree
$2$ is \Cref{thm:lef11}, and degree $2n-2$ follows from it. For $n\le3$
every even degree is one of these.
\end{proof}

\subsection{Surjective morphisms and blow-ups}

\begin{proposition}\label{prop:descent}
Let $f\colon Z\to X$ be a surjective morphism, with $X$ connected, and let
$r=\dim Z-\dim X$. Let $h_{Z}$ be an ample class on $Z$, and $d>0$ the degree
of $h_{Z}^{r}$ on a general fibre. Then $\alpha=d^{-1}f_{*}(h_{Z}^{r}\cup
f^{*}\alpha)$ for every $\alpha\in H^{*}(X,\QQ)$. In particular the
correspondence $\Gamma=d^{-1}(h_{Z}^{r}\times1)\cdot\Gamma_{f}$ from $Z$ to
$X$ satisfies $\Gamma_{*}f^{*}=\mathrm{id}$, so $\Gamma_{*}$ is surjective,
and \HC{} for $Z$ implies \HC{} for $X$.
\end{proposition}

\begin{proof}
By the projection formula
$f_{*}(h_{Z}^{r}\cup f^{*}\alpha)=f_{*}(h_{Z}^{r})\cup\alpha$, and
$f_{*}(h_{Z}^{r})\in H^{0}(X,\QQ)=\QQ$ is the degree $d$, which is positive
because $h_{Z}$ is ample on the fibres. The last claim follows from
\Cref{lem:dominated}.
\end{proof}

\begin{proposition}\label{prop:blowup}
Let $Y\subseteq X$ be a smooth closed subvariety, possibly disconnected, of
codimension $c\ge2$ in every component, let $\pi\colon\tilde X\to X$ be the
blow-up along $Y$, let $j\colon E\to\tilde X$ be the exceptional divisor and
$q\colon E\to Y$ the projective bundle, and let $\xi=c_{1}(\cO_{E}(1))$. Then
\[
  H^{k}(\tilde X,\QQ)=\pi^{*}H^{k}(X,\QQ)\;\oplus\;
  \bigoplus_{i=0}^{c-2}j_{*}\bigl(\xi^{i}\cup q^{*}H^{k-2i-2}(Y,\QQ)\bigr),
\]
and each summand is the image of a correspondence: $\pi^{*}$ from $X$, and
$y\mapsto j_{*}(\xi^{i}\cup q^{*}y)$ from $Y$.
\end{proposition}

\begin{proof}
The decomposition is the blow-up formula \cite[Theorem~7.31]{Voi02b}. The
map $y\mapsto j_{*}(\xi^{i}\cup q^{*}y)$ is the correspondence given by the
image in $Y\times\tilde X$ of a cycle on $E$ representing $\xi^{i}$, pushed
forward along $(q,j)\colon E\to Y\times\tilde X$.
\end{proof}

\subsection{The statements}

\begin{definition}\label{def:statements}
Let $X$ range over all varieties and $p$ over all integers.
\begin{enumerate}[label=\textup{(\arabic*)}]
\item \st{F2}: the Hodge conjecture for every complex abelian
variety.\footnote{In \cite{BhH8} the name \st{F2} denotes a statement about
the Hodge classes not generated by divisors and Weil classes, and it is
proved there to be equivalent to the Hodge conjecture for abelian varieties.
We use the simpler form.}
\item \st{F3$'$}: $\Hdg^{p}(X)=\Ab^{p}(X)$, where $\Ab^{p}(X)$ is the span of
the classes $\Gamma_{*}\gamma$ of degree $2p$ with $B$ an abelian variety
(the point included), $\gamma$ a Hodge class on $B$ and $\Gamma$ an
algebraic cycle on $B\times X$ \cite[Definition~20.34]{BhH8}.
\item \st{L}: for every $X$ and every ample class on $X$, the Lefschetz
involution $*_{L}$ defined below is induced by an algebraic cycle on
$X\times X$. This is Grothendieck's standard conjecture $B(X)$ for every $X$
\cite{Gro69}: by \cite[1.4.4]{Kle68} and \cite[Proposition~1.2]{And96}, the
$\QQ$-algebra of operators generated by $L$ and $*_{L}$ is the one generated by
$L$ and the operator $\Lambda$ of $B(X)$.
\item \st{M}: every Hodge class is motivated in the sense of
\cite[\S1]{And96}.
\item \st{V}: for every smooth projective morphism $f\colon\cX\to S$ over a
smooth connected quasi-projective variety, and every global section $\xi$ of
$R^{2p}f_{*}\QQ$ over $S^{\mathrm{an}}$, if $\xi_{s}$ is algebraic for one
$s\in S$ then $\xi_{t}$ is algebraic for every $t\in S$.
\item \st{CM}: the Hodge conjecture for every abelian variety of CM type.
\item \st{IP}: the case of \st{V} in which $f$ is a polarized abelian scheme
and $\cX_{s}$ is of CM type.
\end{enumerate}
\end{definition}

An abelian variety $A$ is \emph{of CM type} if $\End(A)\otimes\QQ$ contains a
commutative semisimple $\QQ$-algebra of dimension $2\dim A$; the point is of
CM type, and products of abelian varieties of CM type are of CM type.

Each $\Gamma_{*}\gamma$ in \st{F3$'$} is a Hodge class, so
$\Ab^{p}(X)\subseteq\Hdg^{p}(X)$; taking $B$ a point shows that $\Ab^{p}(X)$
contains $\Alg^{p}(X)$. Thus \st{F3$'$} asks only that Hodge classes be
algebraic \emph{modulo} what abelian varieties supply. By
\Cref{prop:descent}, applied to $A\times\PP^{1}\to A$, the conjecture for
varieties that are not abelian already contains the conjecture for abelian
varieties; this is why the half of the conjecture that concerns other
varieties has to be stated modulo abelian varieties to be a separate
problem.

\subsection{The Lefschetz standard conjecture and motivated classes}

Let $X$ have dimension $n$ and ample class $h$. Every class in $H^{k}(X)$ is
uniquely $\sum_{j}L^{j}x_{j}$ with $x_{j}$ primitive of degree $k-2j$, and the
\emph{Lefschetz involution} is defined by $*_{L}(L^{j}x_{j})=L^{n-k+j}x_{j}$
\cite[\S1.1]{And96}. On $H^{k}$ with $k\le n$ it is $L^{n-k}$, and on
$H^{2n-k}$ it is the inverse of $L^{n-k}$. A class $\alpha\in H^{2p}(X,\QQ)$
is \emph{motivated} \cite[D\'efinition~1]{And96} if there are a variety $Y$,
ample classes $h_{X}$ on $X$ and $h_{Y}$ on $Y$, and algebraic classes
$\beta,\gamma$ on $X\times Y$ with
\[
  \alpha=\pr_{X*}\bigl(\beta\cup *_{L}\gamma\bigr),
\]
where $*_{L}$ is the Lefschetz involution of $X\times Y$ for the ample class
$\pr_{X}^{*}h_{X}+\pr_{Y}^{*}h_{Y}$.
We use three facts from \cite{And96}. (a) Motivated classes are Hodge
classes, they contain the algebraic classes, and they form a graded
$\QQ$-algebra that is stable under pull-back and push-forward along
projections \cite[Proposition~2.1 and \S2.5]{And96}; hence a correspondence
given by an algebraic cycle maps motivated classes to motivated
classes.\footnote{For a correspondence, $\Gamma_{*}\gamma=\pr_{X*}(\cl
(\Gamma)\cup\pr_{B}^{*}\gamma)$ is built from $\gamma$ by a pull-back along a
projection, a product with an algebraic class and a push-forward along a
projection.} (b) Every Hodge class on an abelian variety is motivated
\cite[Th\'eor\`eme~0.6.2]{And96}. (c) If the Lefschetz involution of every
variety, for every ample class, is induced by an algebraic cycle, then every
motivated class is algebraic \cite[\S2, after D\'efinition~1]{And96}. Fact
(c) is immediate from the definition: $*_{L}\gamma$ is then the image of an
algebraic class under an algebraic correspondence.

\begin{proposition}\label{prop:LM}
\st{L} and \st{M} together imply \HC.
\end{proposition}

\begin{proof}
By \st{M} every Hodge class is motivated, and by \st{L} and fact~(c) every
motivated class is algebraic.
\end{proof}

\begin{proposition}\label{prop:f3m}
\st{F3$'$} implies \st{M}.
\end{proposition}

\begin{proof}
Under \st{F3$'$} every Hodge class is a sum of classes $\Gamma_{*}\gamma$ with
$\gamma$ a Hodge class on an abelian variety. Such a $\gamma$ is motivated by
fact~(b), and $\Gamma_{*}\gamma$ is motivated by fact~(a).
\end{proof}

\subsection{Families}

\begin{lemma}[Theorem of the fixed part]\label{lem:fixed}
Let $f\colon\cX\to S$ be a smooth projective morphism over a smooth connected
quasi-projective variety, and $\xi$ a global section of $R^{k}f_{*}\QQ$. If
$\xi_{s}$ is a Hodge class for one $s\in S$, then $\xi_{t}$ is a Hodge class
for every $t\in S$.
\end{lemma}

\begin{proof}
By Deligne's theorem of the fixed part \cite[Corollaire~4.1.2]{Del71} the
space $H^{0}(S,R^{k}f_{*}\QQ)$ carries a Hodge structure for which each
restriction $r_{t}\colon H^{0}(S,R^{k}f_{*}\QQ)\to H^{k}(\cX_{t},\QQ)$ is an
injective morphism of Hodge structures.\footnote{Injective because a section
of a local system over a connected base is determined by its value at one
point. The Hodge structure is the image of $H^{k}(\bar\cX,\QQ)\to
H^{k}(\cX_{t},\QQ)$ for a smooth compactification $\bar\cX$ of $\cX$, which
by \cite[4.1.1]{Del71} is the whole invariant part. A global section that is
of type $(p,q)$ at one point is of type $(p,q)$ everywhere
\cite[4.1.3.1]{Del71}.} An injective morphism
of Hodge structures that sends a class $x$ to a class of type $(p,p)$
forces $x$ to be of type $(p,p)$, because it respects the Hodge
decompositions. So $\xi$ is a Hodge class of $H^{0}(S,R^{k}f_{*}\QQ)$, and
$\xi_{t}=r_{t}(\xi)$ is a Hodge class for every $t$.
\end{proof}

\begin{lemma}\label{lem:hcimplies}
\HC{} implies each of \st{F2}, \st{F3$'$}, \st{L}, \st{M}, \st{V}, \st{CM}
and \st{IP}.
\end{lemma}

\begin{proof}
\st{F2} and \st{CM} are special cases, and \st{F3$'$} holds because
$\Ab^{p}(X)$ contains $\Alg^{p}(X)$. For \st{L}: on each $H^{k}(X)$ the involution $*_{L}$ is
a morphism of Hodge structures $H^{k}(X)\to H^{2n-k}(X)(n-k)$, so under the
K\"unneth decomposition and Poincar\'e duality each of its components is a
Hodge class on $X\times X$,\footnote{A $\QQ$-linear map $u\colon H^{a}(X)\to H^{b}(X)$
corresponds to an element of $H^{a}(X)^{\vee}\otimes H^{b}(X)\cong
H^{2n-a}(X)(n)\otimes H^{b}(X)$, a summand of $H^{2n-a+b}(X\times X)(n)$,
and the correspondence it defines acts as $u$. The map $u$ is a morphism of
Hodge structures $H^{a}(X)\to H^{b}(X)(r)$ exactly when this element is a
Hodge class of the appropriate weight.} and \HC{} makes it algebraic.
Algebraic classes are motivated, which gives \st{M}. For \st{V}: if $\xi_{s}$
is algebraic, it is a Hodge class, so $\xi_{t}$ is a Hodge class for every
$t$ by \Cref{lem:fixed}, and \HC{} makes it algebraic. \st{IP} is a case of
\st{V}.
\end{proof}
